\documentclass{article}

% packages
\usepackage{amsmath, amsthm, thmtools, amsfonts, amssymb, luacode, catchfile, tikzducks, hyperref, ifthen}
\ifcsname c@kobocompile\endcsname
	\usepackage[a5paper, total={1072pt, 1448pt}, margin=10pt, includeheadfoot]{geometry} % set page margins
\else
	\usepackage[a4paper, margin=50pt, includeheadfoot]{geometry}
\fi
\usepackage[shortlabels]{enumitem}
\usepackage[skip=3pt, indent=0pt]{parskip}

% language
\usepackage[bidi=basic, layout=tabular, provide=*]{babel}
\ifcsname c@english\endcsname
	\babelprovide[main, import]{english}
\else
	\babelprovide[main, import]{hebrew}
	\babelprovide{rl}
\fi
%\babelfont{rm}{Libertinus Serif}
\babelfont{rm}[Renderer=Harfbuzz]{Libertinus Serif}
\babelfont{sf}{Libertinus Sans}
\babelfont{tt}{Libertinus Mono}

% style
\AddToHook{cmd/section/before}{\clearpage}	% Add line break before section
\linespread{1.3}
\setcounter{secnumdepth}{0}		% Remove default number tags from sections, this won't do well with theorems
\AtBeginDocument{\setlength{\belowdisplayskip}{3pt}}
\AtBeginDocument{\setlength{\abovedisplayskip}{3pt}}
\graphicspath{ {../images/} }

% operators
\DeclareMathOperator\cis{cis}
\DeclareMathOperator\Sp{Sp}
\DeclareMathOperator\tr{tr}
\DeclareMathOperator\im{Im}
\DeclareMathOperator\re{Re}
\DeclareMathOperator\diag{diag}
\DeclareMathOperator*\lowlim{\underline{lim}}
\DeclareMathOperator*\uplim{\overline{lim}}
\DeclareMathOperator\rng{rng}
\DeclareMathOperator\Sym{Sym}
\DeclareMathOperator\Arg{Arg}
\DeclareMathOperator\Log{Log}
\DeclareMathOperator\dom{dom}
\DeclareMathOperator\supp{Supp}
\DeclareMathOperator\var{Var}
\DeclareMathOperator\cov{Cov}

% commands
%\renewcommand\qedsymbol{\textbf{מש''ל}}
%\renewcommand\qedsymbol{\fbox{\emoji{lizard}}}
\newcommand{\Aa}[0]{\mathcal{A}}
\newcommand{\Bb}[0]{\mathcal{B}}
\newcommand{\CC}[0]{\mathbb{C}}
\newcommand{\Cc}[0]{\mathcal{C}}
\newcommand{\EE}[0]{\mathbb{E}}
\newcommand{\FF}[0]{\mathbb{F}}
\newcommand{\Ff}[0]{\mathcal{F}}
\newcommand{\Ii}[0]{\mathcal{I}}
\newcommand{\Gg}[0]{\mathcal{G}}
\newcommand{\Ll}[0]{\mathcal{L}}
\newcommand{\Mm}[0]{\mathcal{M}}
\newcommand{\NN}[0]{\mathbb{N}}
\newcommand{\Nn}[0]{\mathcal{N}}
\newcommand{\PP}[0]{\mathbb{P}}
\newcommand{\Pp}[0]{\mathcal{P}}
\newcommand{\QQ}[0]{\mathbb{Q}}
\newcommand{\RR}[0]{\mathbb{R}}
\newcommand{\Rr}[0]{\mathcal{R}}
\newcommand{\Ss}[0]{\mathcal{S}}
\newcommand{\TT}[0]{\mathbb{T}}
\newcommand{\Uu}[0]{\mathcal{U}}
\newcommand{\Vv}[0]{\mathcal{V}}
\newcommand{\Ww}[0]{\mathcal{W}}
\newcommand{\ZZ}[0]{\mathbb{Z}}
\newcommand{\acts}[0]{\circlearrowright}
\newcommand{\explain}[2] {
	\begin{flalign*}
		 && \text{#2} && \text{#1}
	\end{flalign*}
}
\newcommand{\maketitleprint}[0]{ \begin{center}
	%\begin{tikzpicture}[scale=3]
	%	\duck[graduate=gray!20!black, tassel=red!70!black]
	%\end{tikzpicture}	
	\includegraphics[width=6cm]{cover}
\end{center}
}

% theorem commands
\newtheoremstyle{c_remark}
	{}	% Space above
	{}	% Space below
	{}% Body font
	{}	% Indent amount
	{\bfseries}	% Theorem head font
	{}	% Punctuation after theorem head
	{.5em}	% Space after theorem head
	{\thmname{#1}\thmnumber{ #2}\thmnote{ \normalfont{\text{(#3)}}}}	% head content
\newtheoremstyle{c_definition}
	{3pt}	% Space above
	{3pt}	% Space below
	{}% Body font
	{}	% Indent amount
	{\bfseries}	% Theorem head font
	{}	% Punctuation after theorem head
	{.5em}	% Space after theorem head
	{\thmname{#1}\thmnumber{ #2}\thmnote{ \normalfont{\text{(#3)}}}}	% head content
\newtheoremstyle{c_plain}
	{3pt}	% Space above
	{3pt}	% Space below
	{\itshape}% Body font
	{}	% Indent amount
	{\bfseries}	% Theorem head font
	{}	% Punctuation after theorem head
	{.5em}	% Space after theorem head
	{\thmname{#1}\thmnumber{ #2}\thmnote{ \text{(#3)}}}	% head content

\ifcsname c@english\endcsname
	\theoremstyle{plain}
	\newtheorem{theorem}{Theorem}[section]
	\newtheorem{lemma}[theorem]{Lemma}
	\newtheorem{proposition}[theorem]{Proposition}
	\newtheorem*{proposition*}{Proposition}
	%\newtheorem{corollary}[theorem]{אין חלופה עברית}

	\theoremstyle{definition}
	\newtheorem{definition}[theorem]{Definition}
	\newtheorem*{definition*}{Definition}
	\newtheorem{example}{Example}[section]
	\newtheorem{exercise}{Exercise}[section]

	\theoremstyle{remark}
	\newtheorem*{remark}{Remark}
	\newtheorem*{solution}{Solution}
	\newtheorem{conclusion}[theorem]{Conclusion}
	\newtheorem{notation}[theorem]{Notation}
\else
	\theoremstyle{c_plain}
	\newtheorem{theorem}{משפט}[section]
	\newtheorem{lemma}[theorem]{למה}
	\newtheorem{proposition}[theorem]{טענה}
	\newtheorem*{proposition*}{טענה}
	%\newtheorem{corollary}[theorem]{אין חלופה עברית}

	\theoremstyle{c_definition}
	\newtheorem{definition}[theorem]{הגדרה}
	\newtheorem*{definition*}{הגדרה}
	\newtheorem{example}{דוגמה}[section]
	\newtheorem{exercise}{תרגיל}[section]

	\theoremstyle{c_remark}
	\newtheorem*{remark}{הערה}
	\newtheorem*{solution}{פתרון}
	\newtheorem{conclusion}[theorem]{מסקנה}
	\newtheorem{notation}[theorem]{סימון}
\fi

% Questions related commands
\newcounter{question}
\setcounter{question}{1}
\newcounter{sub_question}
\setcounter{sub_question}{1}

\ifcsname c@english\endcsname
	\newcommand{\question}[1][0]{
		\ifthenelse{#1 = 0}{}{\setcounter{question}{#1}}
		\section{Question \arabic{question}}
		\addtocounter{question}{1}
		\setcounter{sub_question}{1}
	}

	\newcommand{\subquestion}[1][0]{
		\ifthenelse{#1 = 0}{}{\setcounter{sub_question}{#1}}
		\subsection{Part \alph{sub_question}}
		\addtocounter{sub_question}{1}
	}
\else
	\newcommand{\question}[1][0]{
		\ifthenelse{#1 = 0}{}{\setcounter{question}{#1}}
		\section{שאלה \arabic{question}}
		\addtocounter{question}{1}
		\setcounter{sub_question}{1}
	}

	\newcommand{\subquestion}[1][0]{
		\ifthenelse{#1 = 0}{}{\setcounter{sub_question}{#1}}
		\subsection{סעיף \localecounter{letters.gershayim}{sub_question}}
		\addtocounter{sub_question}{1}
	}
\fi

% import lua and start of document
\directlua{common = require ('../common')}

\GetEnv{AUTHOR}

% headers
\author{\AUTHOR}
\date\today

\title{פתרון מטלה 08 --- אנליזה על יריעות, 80426}

\DeclareMathOperator{\vol}{vol}

\begin{document}
\maketitle
\maketitleprint[blue]

\question{}
תהינה $\varphi_1, \ldots, \varphi_k \in C^\infty(\RR^n)$ ו־$p \in \RR^n$ כך ש־$\varphi_1(p) = \cdots = \varphi_k(p) = 0$ וכן $\{ \nabla \varphi_1(p), \ldots, \nabla \varphi_n(p) \}$ בלתי תלויה־לינארית.

\subquestion{}
נראה שקיימת סביבה $U$ של $p$ כך שהקבוצה,
\[
	M
	= \{ x \in U \mid \varphi_1(x) = \cdots = \varphi_k(x) = 0 \}
\]
היא יריעה $(n - k)$־מימדית,
ונחשב את $T_p M$ ב־$p$.
\begin{proof}
	נגדיר פונקציה $f : \RR^n \to \RR^k$ על־ידי $f(x) = (\varphi_1(x), \ldots, \varphi_k(x))$.
	$f$ חלקה ישירות מהגדרתה והנתון אודות $\varphi_i$ לכל $i \le n$.
	אנו יודעים ש־$\dim D f \mid_p = k$ ולכן היא רגולרית וממשפט התמונה ההפוכה $f^{-1}(0)$ היא יריעה.
	נבחין כי זוהי היריעה $N = \{ x \in U \mid \varphi_1(x) = \cdots = \varphi_k(x) \}$, ונוכל לצמצם אותה עם סביבה פתוחה כלשהי $p \in U$ כך ש־$N \cap U$ לא טריוויאלי ולקבל $M$ כזו.
	נעיר כי ישירות מהמשפט זוהי יריעה $(n - k)$־מימדית.
\end{proof}
נעבור לחישוב $T_p M$.
אנו יודעים ממשפט התמונה ההפוכה שמתקיים,
\[
	T_p M
	= \ker df_p
	= \{ x \in \RR^n \mid (\nabla \varphi_1(p), \ldots, \nabla \varphi_k(p)) \cdot x = 0 \}
	= \RR^n \setminus \Sp\{ \varphi_1(p), \ldots, \varphi_k(p) \}
\]

\subquestion{}
תהי $f : \RR^n \to \RR$ העתקה חלקה כך ש־$f(x) \ge f(p)$ לכל $x \in M$.
נראה ש־$\nabla f(p) \perp T_p M$.
\begin{proof}
	נבחין כי $D_u f \ge 0$ לכל $u$ וקטור כיוון ב־$M$, ולכן נובע בפרט ש־$D_u f = 0$ עבור $u \in T_p M$ ישירות מהגדרה, כלומר נובע ש־$\nabla f(p) \perp T_p M$ כפי שרצינו.
\end{proof}

\subquestion{}
נראה שהמשלים האורתוגונלי של $T_p M$ הוא $\Sp\{ \nabla \varphi_1(p), \ldots, \nabla \varphi_k(p) \}$.
\begin{proof}
	%מתברר שהסעיף טיפשי ממשפט התמונה ההפוכה, מוכיחים מהאפיון מהאפיון של המרחב המשיק כאוסף נגזרות של מסילות בנקודה, ו־$T_p M = \ker df_p$.
	ישירות כמסקנה ממשפט התמונה ההפוכה ואורתוגונליות תמונה וגרעין העתקה לינארית,
	\[
		T_p M
		= \ker\{ \nabla \varphi_1(p), \ldots, \varphi_k(p) \}
		\perp \Sp\{ \nabla \varphi_1(p), \ldots, \varphi_k(p) \}
	\]
	והטענה נובעת ישירות.
\end{proof}

\subquestion{}
נסיק את נוסחת כופלי לגרנז'.
\begin{proof}
	נבחין כי אם $x$ מינימום מקומי של $f$ אז היא במשלים האורתוגונלי של $T_p M$ ולכן בלתי תלויה־לינארית בה.
\end{proof}

\question{}
נראה שכל יריעה היא תמונה הפוכה של העתקה רגולרית,
ונסיק שחיתוך רוחבי של תת־יריעות הוא יריעה.

\subquestion{}
תהי $N$ יריעה ונניח ש־$Z \subseteq N$ תת־יריעה של $N$.
נראה שלכל $p \in Z$ קיימת סביבה פתוחה $p \in U \subseteq N$ ו־$f : U \to \RR^l$ חלקה כך ש־$U \cap Z = f^{-1}(\{ 0 \})$ כאשר $l = \dim N - \dim Z$.
\begin{proof}
	נתחיל בהוכחת הטענה עבור המקרה $N = \RR^n$ ו־$Z^k \subseteq \RR^n$.
	תהי נקודה $p \in M$ ותהי $\alpha : V \to Z$ פרמטריזציה מקומית כך ש־$V \subseteq \RR^k$ פתוחה.
	מקיום הופכי משמאל לפרמטרזיציה מקומית נסיק שקיימת $\beta : W \to \RR^k$ העתקה חלקה, כאשר $p \in W \subseteq \RR^n$ קבוצה פתוחה.
	מתקיים $\beta \circ \alpha = \id$.
	נגדיר $f(x) = \beta(x) - x$ ונקבל ש־$f(x) = 0 \iff x \in \alpha(W \cap Z)$, כלומר $f^{-1}(\{ 0 \}) = W \cap Z$.

	נעבור למקרה הכללי, נניח ש־$Z^k \subseteq N^h \subseteq \RR^n$ יריעה ותת־יריעה.
	נניח ש־$p \in Z$ ולכן בפרט $p \in N$ וקיימת פרמטריזציה $\alpha : V \to N$ כך ש־$V \subseteq \RR^h$ פתוחה.
	נבחן את $Z_0 = \alpha^{-1}(Z \cap \alpha(V))$, זוהי יריעה $k$־מימדית ב־$\RR^h$ ולכן מקיימת את תנאי החלק הראשון, וקיימת $f : U \to \RR^{h - k}$ כך ש־$U \subseteq \RR^h$.
	נובע אם כך ש־$f \circ \alpha : V \to \RR^{h - k}$ העתקה חלקה המקיימת את הטענה, היא חלקה ולכן קיימת לה הרחבה $\tilde{f} : \tilde{V} \to \RR^{h - k}$ חלקה כך ש־$\tilde{V} \subseteq \RR^n$.
	נוכל להגדיר אותה כך ש־$x \notin V \implies \tilde{f}(x) \ne 0$, ונקבל ש־$\tilde{f}$ מקיימת את הטענה.
\end{proof}

\subquestion{}
נסיק שאם $X, Y$ הן תת־יריעות רוחביות של $N$ אז $X \cap Y$ יריעה ממימד $\dim Y + \dim X - \dim N$.
\begin{proof}
	אילו היריעות זרות אז $X \cap Y$ היא יריעה באופן מנוון, ולכן היא $0$־מימדית, ומרוחביות אכן $\dim Y + \dim X = \dim N$.

	נניח ש־$X \cap Y \ne \emptyset$.
	נגדיר $f : X \to N$ חלקה על־ידי שיכון.
	נבחין כי $f^{1}(Y) \ne \emptyset$, ולכן $f^{-1}(Y) = X \cap Y$ היא יריעה $(\dim Y + \dim X - \dim N)$־מימדית ממשפט מהתרגול.
\end{proof}

\question{}
\subquestion{}
נניח ש־$f : \RR^n \to \RR$ העתקה חלקה ונניח ש־$a \in \RR$ ערך רגולרי של $f$ כך ש־$f^{-1}(\{ a \}) \ne \emptyset$.
נראה ש־$M = \{ x \in \RR^n \mid f(x) \ge a \}$ היא יריעה $n$־מימדית עם השפה $f^{-1}(\{ a \})$.
\begin{proof}
	נבחין כי $\{ x \in \RR^n \mid f(x) > a \}$ היא יריעה $n$־מימדית כמקור של קבוצה פתוחה של פונקציה חלקה (בפרט רציפה).
	לכן מספיק שנראה שלכל $p \in f^{-1}(\{ a \})$ יש פרמטריזציה מקומית (עם שפה) ונוכל להסיק ש־$M$ היא יריעה $n$־מימדית. \\
	תהי $p$ כזו, אז $p$ היא נקודה רגולרית של $f$ מהגדרת ערך רגולרי, כלומר $D f |_p$ הוא על $\RR^n$.
	אז ניקח את ההעתקה לא יודע

	נעיר כי בהוכחה התבססנו על הטענה כי כל יריעה ללא שפה היא יריעה עם שפה.
\end{proof}

\end{document}
